\subsection{Poincaré group of topological groups}

If G is a topological group, for every \(g \in G\) , the left translation \(x \mapsto gx\) is a homeomorphism of G onto itself (TG, III, p.~2) and it induces an isomorphism of \(\pi_{1}(G, e)\) onto \(\pi_{1}(G, g)\) .

Let G be a topological group, let e be its identity element, and denote by \(G_{0}\) the path-connected component of e in G. Denote by R the equivalence relation on G whose equivalence classes are the path-connected components of G, and let \(p: G \to \pi_{0}(G)\) be the canonical map. Since every translation, left or right, is a homeomorphism of G, the relation R is compatible on the left and on the right with the group law of G. By th. 2 of A, I, p.~35, the group \(G_{0}\) is a normal subgroup of G, the quotient group \(G/G_{0}\) is equal to \(\pi_{0}(G)\) endowed with the composition law induced by that of G, and the map p is a group homomorphism.

The group \(\pi_{1}(G,e)\) is called the Poincaré group of G and is denoted by \(\pi_{1}(G)\) . The canonical injection of \(G_{0}\) into G induces an isomorphism of \(\pi_{1}(G_{0})\) onto \(\pi_{1}(G)\) (III, p.~293, remark 2).

Let \(g\) be an element of G; the right translation \(\boldsymbol{\delta}_g \colon x \mapsto xg\) induces an isomorphism \((\boldsymbol{\delta}_g)_* \colon \pi_1(\mathrm{G}) \to \pi_1(\mathrm{G}, g)\) . Suppose that \(g\) belongs to \(\mathrm{G}_0\) and let \(b\) be a path connecting \(e\) to \(g\) . The map \(\sigma \colon \mathrm{G} \times \mathbf{I} \to \mathrm{G}\) defined by \(\sigma(x, t) = xb(t)\) is a homotopy connecting \(\mathrm{Id}_{\mathrm{G}}\) to \(\boldsymbol{\delta}_g\) . By prop. 3 of III, p.~295, one therefore has, for every \(\alpha \in \pi_1(\mathrm{G})\) , \((\boldsymbol{\delta}_g)_*(\alpha) = \beta^{-1}\alpha\beta\) , where \(\beta \in \varpi_{e,g}(\mathrm{G})\) is the class of the path \(b\) . If one denotes by \(\gamma_g\) the left translation, \(x \mapsto gx\) , one likewise has \((\gamma_g)_*(\alpha) = \beta^{-1}\alpha\beta\) .

For every \(g \in G\) , the map \(\operatorname{Int}(g)\) : \(G \to G\) induces an automorphism \(\pi_{1}(\operatorname{Int}(g))\) of \(\pi_{1}(G)\) . One thus defines an action of G on \(\pi_{1}(G)\) . For all \(g \in G\) , we have \(\pi_{1}(\operatorname{Int}(g)) = (\boldsymbol{\gamma}_{g^{-1}})_{*} \circ (\boldsymbol{\delta}_{g})_{*}\) , so that the subgroup \(G_{0}\) operates trivially; this yields an action of \(\pi_{0}(G)\) on \(\pi_{1}(G)\) . When G is a commutative group, this action is trivial, but this is not always the case (IV, p.~459, exerc. 1).

